\lesson{Explicit evidence for subcapturing}\label{sec:capture}
\focus{Subcapturing evidence relates capture sets through their interpretation as capability roots.}

The core argument checked evidence under assumptions, then closed it over
a typed store before interpreting it. The capture extension adds another
kind of evidence to the same process. Its interpretation will bound
capability roots; substituting a call's argument into that evidence will
connect a closure's advertised captures to the uses of its executing body.

A capturing type $S^C$ pairs a value's shape $S$ with a capture set $C$
describing its associated capabilities. In the examples,
$\ty{Unit}=\top^\varnothing$ abbreviates the pure universal shape, and
$\ev{unit}$ is a stored pure value of that type. This abbreviation uses
the existing type language without adding a primitive unit type.

Assume the platform supplies two distinct rigid capability roots,
$\kappa_{\ell}$ and $\kappa_c$. A rigid root denotes itself; a program
name is interpreted through its capture annotation. The store contains
closures $\ev{log}$ and $\ev{console}$ with annotations
$\{\kappa_{\ell}\}$ and $\{\kappa_c\}$, respectively. Both are identity
closures on $\ty{Unit}$, with nonempty annotations that serve as
conservative allowances. Their names indicate the intended capability
roles; logging and console I/O are outside the operational semantics.

\subsection*{Subcapturing and root interpretation}
The sets $\{\ev{log}\}$ and $\{\kappa_{\ell}\}$ contain different names,
so they are unrelated by ordinary subset inclusion. Subcapturing also
uses the annotation of $\ev{log}$, which accounts for its capabilities
through $\kappa_{\ell}$:
\[
 \ev{capvar}(\ev{log}):
       \{\ev{log}\}\le_c\{\kappa_{\ell}\}.
\]
In general, if a program name $x$ has type $S^C$, then
$\ev{capvar}(x)$ proves $\{x\}\le_c C$.
The constructor $\ev{elem}$ represents ordinary subset inclusion:
if every element of $C$ is listed in $D$, then
$\ev{elem}(C,D):C\le_c D$.
Transitivity composes capture evidence as it does type evidence.
Union combines evidence for two sets with a shared destination:
from $p:C\le_c D$ and $q:C'\le_c D$, the expression
$\ev{union}(p,q)$ proves $C\cup C'\le_c D$.

Capture evidence is interpreted through \emph{roots}. For the acyclic
annotations used here, $\roots{C}$ follows the capture information
associated with the names in $C$ to their underlying capability roots:
\[
 \roots{\{\ev{log}\}}=\{\kappa_{\ell}\},\qquad
 \roots{\{\ev{console}\}}=\{\kappa_c\},\qquad
 \roots{\{\ev{unit}\}}=\varnothing.
\]
Thus $\roots{C}$ is a set of capability roots, while $\ev{root}(a)$ is
the single program variable obtained by stripping an atom's wrappers.

Over a typed store, well-typed capture evidence has the interpretation
\[
  e:C\le_c D\quad\Longrightarrow\quad
  \roots{C}\subseteq\roots{D}.
\]
This canonical-forms result depends on the typed-store hypothesis.
As with type members, capture members can supply assumptions whose
semantic interpretation requires an actual stored value.

\subsection*{Immediate use sets}
The use-set function assigns a syntactic use bound to each term.
The cases needed for these examples are:
\[
\tgt{\begin{aligned}
\uses{a}&=\{\ev{root}(a)\},& \uses{v}&=\varnothing,\\
 \uses{a\,b}&=\{\ev{root}(a),\ev{root}(b)\},&
 \uses{\ev{let}\ x=t\ \ev{in}\ u\langle U;q\rangle}
   &=\uses{t}\cup U.
\end{aligned}}
\]
Here $v$ denotes a value expression, such as a lambda literal, whose
creation leaves its body unevaluated. A stored name $x$ is an atom with
use set $\{x\}$; its roots are determined by its annotation. Thus the
returned name $\ev{unit}$ has a singleton use set and no capability
roots. The term \emph{immediate} excludes the latent bodies of value
expressions while still including pending let bodies.

In the let's annotation $\langle U;q\rangle$, $U$ is an upper bound
for the pending body's uses, expressed using names available outside
the let. In the body's context, where $x$ is available, $q$ proves
$\uses{u}\le_c U$. Replacing a local name by outer names that account
for it is called \emph{avoidance}. The factory example below includes
an explicit construction of this evidence.

An atom's wrappers affect its typing but preserve its root. In particular,
$\uses{\cast{a}{e}}=\{\ev{root}(a)\}$ when the cast is an atom,
and a term cast satisfies $\uses{\cast{t}{e}}=\uses{t}$.
Evidence can mention names without adding them to immediate uses:
this function bounds the program's modeled inspections, not all names
appearing in its syntax. These clauses matter when an argument is a
wrapped atom rather than a bare variable.

\subsection*{Translating closing evidence}
The source judgment $\src{C;\Gamma\vdash t:T}$ assigns type $T$ to $t$
with permitted immediate uses $C$. A source lambda contains its parameter
type and body; its typing derivation selects the capture annotation $A$
and checks the body with allowance $A\cup\{u\}$. Translation makes this
annotation and its closing evidence explicit in the target lambda.

In Figure~\ref{fig:capture-translation}, translation brackets also apply
to capture sets, types, and contexts. The source context $\Gamma$ is
assumed well formed, and $\mathrm{wf}$ denotes well-formedness.
The judgment symbol $\vdash^c$ denotes capture-evidence typing. The superscript
$\uparrow$ preserves references to outer bindings under a new parameter
binder. With distinct printed names, those names keep their spelling;
the named calculations below leave this shift implicit.

The translation computes two components from the body derivation: a
target body and evidence that its computed use set lies within the source
allowance. The target lambda stores both components, just as the
translation of source subsumption records inclusion evidence in a cast.


\begin{figure}[!htbp]
\begin{examplepanel}{srccol}{Source lambda rule}
\small
\[
\src{\inferrule*[right=\textsc{All-I}]{
 d_b:A^\uparrow\cup\{u\};\Gamma,u:T\vdash t:T'\\
 T\ \mathrm{wf}
}{\varnothing;\Gamma\vdash\lambda(u:T).t:((u:T)\to T')^A}}
\]
Only the parameter type and body occur in the source lambda. Its typing
derivation selects the capture annotation and records the body's permitted uses.
\end{examplepanel}
\begin{examplepanel}{tgtcol}{Target lambda with closing evidence}
\[
\tgt{
 \lambda^{\tr{A}}(u:\tr{T}).\
       \underbrace{\tr{d_b}}_{\text{body}}\
       \left\langle
       \underbrace{\tr{d_b}_{\mathrm{uses}}}_{\text{closing evidence }g}
       \right\rangle .}
\]
\small
\[
\tgt{
\tr{\Gamma},u:\tr{T}\vdash^c
 g:\uses{\tr{d_b}}\le_c\tr{A}^{\uparrow}\cup\{u\}.}
\]
\end{examplepanel}
\caption{Translation of a lambda premise in a well-formed source context
\(\Gamma\). The two components are the translated body \(\tr{d_b}\)
and its closing evidence \(\tr{d_b}_{\mathrm{uses}}\).
The superscript \(\uparrow\) carries outer names under the new parameter
binder. Angle brackets delimit stored evidence; the lambda receives only
its term argument at runtime. The functions \texttt{translate} and
\texttt{translateUses} construct these components.}
\label{fig:capture-translation}
\end{figure}


\FloatBarrier
\subsection*{Computing a closing certificate from a derivation}
The translation contract in Figure~\ref{fig:capture-translation} tells us
what $g$ proves. To see how it is built, translate the source body
$\ev{log}\ u$ rather than choosing its certificate by hand.
Write $F=\ty{Unit}\to\ty{Unit}$ and $U=\{\ev{log},u\}$.
All the following judgments are under
$\Delta=\Gamma,u:\ty{Unit}$, where
$\Gamma(\ev{log})=F^{\{\kappa_\ell\}}$.

The source variable rule records the variable itself in both its use set
and its assigned capture set:
\[
 \src{\{\ev{log}\};\Delta\vdash\ev{log}:F^{\{\ev{log}\}}},
 \qquad
 \src{\{u\};\Delta\vdash u:\top^{\{u\}}}.
\]
In particular, the second conclusion is not yet the argument type
$\ty{Unit}=\top^\varnothing$. Source subsumption can change
\emph{both} an assigned type and a permitted use set:
\[
 \src{\frac{d:C;\Delta\vdash t:S\qquad e:\Delta\vdash S<:T
                    \qquad f:\Delta\vdash C\le_c D}
           {\textsc{Sub}(d,e,f):D;\Delta\vdash t:T}}.
\]
These are independent obligations. For the type change on $u$, its
declared pure type justifies $\{u\}\le_c\varnothing$ by the source
variable subcapturing rule. For its use allowance, ordinary inclusion
justifies $\{u\}\le_c U$. For $\ev{log}$, we similarly recover its
declared type using $\{\ev{log}\}\le_c\{\kappa_\ell\}$, and widen its
permitted uses using $\{\ev{log}\}\le_c U$.

Let $d_\ell$ and $d_u$ be these two subsumed variable derivations.
Both have allowance $U$, so the source application rule combines them:
\[
 \src{\frac{d_\ell:U;\Delta\vdash\ev{log}:F^{\{\kappa_\ell\}}
                   \qquad d_u:U;\Delta\vdash u:\ty{Unit}}
           {d_b=\textsc{App}(d_\ell,d_u):
                      U;\Delta\vdash\ev{log}\ u:\ty{Unit}}}.
\]
This specifies the derivation whose translation we will compute.

\textbf{The body component.}
Two target constructors record the capture information just used.
$\ev{recap}(a,p)$ preserves an atom's shape and assigns capture set
$D$ when $p:\{\ev{root}(a)\}\le_c D$.
$\ev{capt}(e,p)$ combines shape inclusion $e$ with capture inclusion
$p$ to convert a capturing type. Thus variable translation and the two
type-subsumption premises give the exact atoms
\[
\tgt{\begin{aligned}
 L&=\cast{\ev{recap}(\ev{log},\ev{refl}\{\ev{log}\})}
          {\ev{capt}(\ev{refl}(F),\ev{capvar}(\ev{log}))},\\
 B&=\cast{\ev{recap}(u,\ev{refl}\{u\})}
          {\ev{capt}(\ev{refl}(\top),\ev{capvar}(u))}.
\end{aligned}}
\]
Here reflexivity on a capture set is capture evidence.
The $\ev{capvar}$ expressions read the declared types of the bare
bindings. They are the translations of the source variable-subcapturing
premises. The resulting body is $\tr{d_b}=L\,B$, with
$\uses{L\,B}=\{\ev{log},u\}$ because both wrappers preserve roots.

\textbf{The use-evidence component.}
Only three clauses of $\tr{-}_{\mathrm{uses}}$ are needed:
\[
\tgt{\begin{aligned}
 \tr{\textsc{Var}(x)}_{\mathrm{uses}}&=\ev{refl}\{x\},\\
 \tr{\textsc{Sub}(d,e,f)}_{\mathrm{uses}}
     &=\tr{d}_{\mathrm{uses}};\tr{f},\\
 \tr{\textsc{App}(d_1,d_2)}_{\mathrm{uses}}
     &=\ev{union}(\tr{d_1}_{\mathrm{uses}},
                  \tr{d_2}_{\mathrm{uses}}).
\end{aligned}}
\]
The type-subtyping premise $e$ affects the body cast; the distinct
use-subcapturing premise $f$ affects this certificate.
Applying the clauses from the variable leaves upwards gives
\[
\tgt{\begin{aligned}
 p_\ell&=\ev{refl}\{\ev{log}\};
                  \ev{elem}(\{\ev{log}\},U),\\
 p_u&=\ev{refl}\{u\};\ev{elem}(\{u\},U),\\
 g_{\mathrm{gen}}&=\ev{union}(p_\ell,p_u):
                     \{\ev{log},u\}\le_c U.
\end{aligned}}
\]
The two branches share destination $U$, exactly as required by
$\ev{union}$. Finally, lambda translation stores both computed components:
\[
 \tgt{\lambda^{\{\ev{log}\}}(u:\ty{Unit}).
                 L\,B\ \langle g_{\mathrm{gen}}\rangle.}
\]
This is the generated syntax for the specified source derivation, in
named notation. The simpler certificate $\ev{elem}(U,U)$ has the same
endpoints but is a different expression. We use that hand-written
certificate in the next trace to keep the operational calculation short.
The substitution argument after the trace applies to either certificate.


\needspace{14\baselineskip}
\subsection*{Closing a closure body}
To isolate the execution argument, consider the simpler hand-written
target closure
\[
 h=\lambda^{\{\ev{log}\}}(u:\ty{Unit}).\,\ev{log}\ u.
\]
Its superscript records the capture annotation $\{\ev{log}\}$, while
its body has immediate use set $\{\ev{log},u\}$. Since the argument is
available when the closure is called, the closing obligation is
\[
 g:\uses{\ev{log}\ u}\le_c\{\ev{log}\}\cup\{u\}.
\]
The two endpoints are identical, so $\ev{elem}$ supplies $g$.
This shares the generated closure's use-set obligation, while its body
and certificate omit the translation's additional syntax.

Figure~\ref{fig:capture-call} gives the closing certificate, root
calculation, and complete two-call trace for this hand-written target
example. A body produced by source translation may contain additional
atom wrappers, with generated closing evidence built from unions and
compositions.


\begin{figure}[!htbp]
\begin{examplepanel}{tgtcol}{Stored closure and closing certificate}
\[
\tgt{\begin{aligned}
 h&=\lambda^{\{\ev{log}\}}(u:\ty{Unit}).\
          \ev{log}\ u\ \langle g_h\rangle,\\
 g_h&=\ev{elem}(\{\ev{log},u\},\{\ev{log},u\}),\\
 g_h&:\uses{\ev{log}\ u}\le_c\{\ev{log}\}\cup\{u\}.
\end{aligned}}
\]
The store also contains \(\ev{unit}:\ty{Unit}\) and the identity closure
\(\ev{log}:(\ty{Unit}\to\ty{Unit})^{\{\kappa_\ell\}}\).
\end{examplepanel}
\begin{examplepanel}{runcol}{Capture-annotation resolution}
\[
\run{\{h\}\ \rightsquigarrow\ \{\mathsf{log}\}\
   \rightsquigarrow\ \{\kappa_\ell\},
   \qquad \operatorname{roots}_{\sigma}(\{\mathsf{unit}\})=\varnothing.}
\]
The relation \(\rightsquigarrow\) resolves capture annotations.
\end{examplepanel}
\begin{examplepanel}{runcol}{Execution with an empty continuation}
\small
\renewcommand{\arraystretch}{1.3}
\begin{tabularx}{\linewidth}{@{}lXll@{}}
\toprule
Current term & Use set & Its roots & Inspection\\
\midrule
$h\ \mathsf{unit}$ & $\{h,\mathsf{unit}\}$ & $\{\kappa_\ell\}$ & $h$\\
$\mathsf{log}\ \mathsf{unit}$ & $\{\mathsf{log},\mathsf{unit}\}$ & $\{\kappa_\ell\}$ & $\mathsf{log}$\\
$\mathsf{unit}$ & $\{\mathsf{unit}\}$ & $\varnothing$ & none\\
\bottomrule
\end{tabularx}
\end{examplepanel}
\caption{Closing evidence, root resolution, and execution of the closure.
The two applications inspect \(h\) and \(\mathsf{log}\), then return an
atom with a singleton use set and no capability roots.}
\label{fig:capture-call}
\end{figure}



Each application inspects the closure named in the last column: the
first executes $h$'s body and the second executes the identity body of
$\ev{log}$, returning the atom $\ev{unit}$. The syntactic use sets change
along this execution while their roots remain bounded by the initial root set.

\FloatBarrier
\needspace{12\baselineskip}
\subsection*{Capture prediction for application}
\label{sec:app-prediction}
We now derive the application bound for an arbitrary argument atom,
including its wrappers. Fix a typed store $\sigma:\Gamma_\sigma$ with
$\sigma(h)=\lambda^A(u:T).t\langle g\rangle$, and suppose
$\Gamma_\sigma\vdash b:T$. The argument must have the stored lambda's
domain type; a coerced call first supplies the required domain cast.
For the present derivation the function receiver is the bare location
$h$, and the continuation is empty.

\textbf{1. Recover the open certificate.}
Typing of the stored lambda gives
\[
 \Gamma_\sigma,u:T\vdash^c
       g:\uses{t}\le_c A^\uparrow\cup\{u\}.
\]
The body and its certificate are still open in $u$. Store typing is what
connects this certificate to the lambda actually found at $h$; an
annotation on a hypothetical parameter alone would not supply that body.

\textbf{2. Substitute the typed argument.}
Put $r=\ev{root}(b)$ and write $\theta=[u\mapsto r]$ for substitution
on capture sets. Term and evidence substitution replace term-level
occurrences of $u$ by the whole atom $b$; their type and capture names
use $r$. The capture-evidence substitution lemma gives
\[
 \Gamma_\sigma\vdash^c g[u:=b]:
       (\uses{t})\theta\le_c(A^\uparrow\cup\{u\})\theta.
\]
This is now evidence over the store alone. Its type still mentions the
substitution of the old use set, which is not yet the displayed use set
of the new body.

\textbf{3. Identify both endpoints.}
Two syntactic identities make the connection:
\[
 \uses{t[u:=b]}=(\uses{t})\theta,\qquad
 (A^\uparrow\cup\{u\})\theta=A\cup\{r\}.
\]
The first is a proved substitution lemma for this particular use-set
function. It is not a consequence of the definition of substitution
alone. For an application, it follows because substitution changes
each operand's root by $\theta$; wrappers add no uses. For a let,
substitution acts on both the right-hand side and its declared allowance
$U$. Value literals have empty immediate uses on both sides.
The second identity removes the parameter binder: $A$ refers only to
outer names, and the singleton parameter becomes $\{r\}$.
Thus the substituted certificate has exactly the endpoints
\[
 \Gamma_\sigma\vdash^c g[u:=b]:
       \uses{t[u:=b]}\le_c A\cup\{r\}.
\]

\textbf{4. Interpret evidence and the stored annotation.}
Capture canonical forms over the typed store turn this judgment into
root inclusion. Separately, store typing gives
$\roots{\{h\}}=\roots{A}$, since $A$ is the annotation of the value
stored at $h$. Combining them gives
\[
 \roots{\uses{t[u:=b]}}
 \subseteq\roots{A\cup\{r\}}
 =\roots{\{h,r\}}
 =\roots{\uses{h\,b}}.
\]
The application step replaces $h\,b$ by $t[u:=b]$, so this is precisely
the required decrease for that step. The inclusion comes from $g$;
the middle equality comes from the stored annotation. Neither is
ordinary subset inclusion between the original sets of program names.

For example, let $b=\cast{\ev{unit}}{\ev{refl}(\ty{Unit})}$ in the
earlier closure. The reduct is $\ev{log}\ b$, not
$\ev{log}\ \ev{unit}$ as target syntax. Nevertheless its use set is
$\{\ev{log},\ev{unit}\}$: the substitution lemma uses
$r=\ev{unit}$. Its roots are $\{\kappa_\ell\}$, exactly as in the
bare-argument trace. This illustrates why the proof keeps the whole
argument in the body and only its root in the capture set.

These are the steps of \texttt{Store.Typed.app\_uses} in the development:
typed substitution, \texttt{Tm.uses\_substAtom}, the closing-set identity,
\texttt{cap\_canon}, and \texttt{Store.Typed.root\_annot}.

\subsection*{Applying the bound through a wrapped receiver}
Now start from a typed application at its \emph{presented} function type:
\[
 \Gamma_\sigma\vdash a:((u:S)\to V)^C,\qquad
 \Gamma_\sigma\vdash b:S.
\]
Let $h=\ev{root}(a)$ and
$\sigma(h)=\lambda^A(u:T).t\langle g\rangle$.
The client's domain $S$ may differ from the stored domain $T$.
After an outer call with argument $b_0$, the generated example has this
form with $a=L$, $b=B[u:=b_0]$, and $h=\ev{log}$.
The use-set clause gives
$\uses{a\,b}=\{h,\ev{root}(b)\}$. We must show that the particular
reduct of the wrapped call has roots below this set.

Canonical forms supply the typing facts needed by the machine's two
wrapped application rules. If the receiver's chain has identity form,
or the capture extension's conversion form preserving the function
shape, the argument already has the stored domain type. The reduct is
$t[u:=b]$, and the preceding four-step argument applies directly.

If the chain instead has function form $\form{pi}(d,c)$, the machine
adapts the argument and result as in Section~\ref{sec:application}.
Write $b'=\cast{b}{d}$ and let $c_b$ be the result evidence instantiated
at the original argument $b$. The reduct is
\[
 t_{\mathrm{next}}=\cast{t[u:=b']}{c_b}.
\]
Typed canonical forms establish $d:S\le T$, hence $b':T$.
The closing-certificate argument therefore applies with $b'$.
The atom-root and term-use clauses then give
\[
 \ev{root}(b')=\ev{root}(b),\qquad
 \uses{t_{\mathrm{next}}}=\uses{t[u:=b']}.
\]
Consequently,
\[
 \roots{\uses{t_{\mathrm{next}}}}
 \subseteq\roots{\{h,\ev{root}(b')\}}
 =\roots{\{h,\ev{root}(b)\}}
 =\roots{\uses{a\,b}}.
\]
The use set records the receiver's root $h$, whose capabilities come
from its stored annotation. The presented capture annotation on $a$'s
type need not be identified with that annotation.

For the generated example's next call, the relevant stored lambda is
$\ev{log}$, so we use its closing certificate. Its identity body has
certificate $\ev{elem}(\{u\},\{\kappa_\ell,u\})$.
The outer certificate $g_{\mathrm{gen}}$ bounded entry into $L\,B$;
the certificate of $\ev{log}$ bounds execution of the subsequently
invoked identity body. Each call uses the evidence belonging to the
closure it actually invokes.

These application steps keep the store and pending continuation $K$
unchanged. A whole state's use set is the union of its current term's
uses with the pending frames' uses. Adding the same $\roots{\uses{K}}$
to both sides therefore gives the application case for a whole state.
The next example shows how those frames acquire their allowances and
why the allowances matter when a continuation is later discharged.

\exercise{roots}{Assume $k$ has capture annotation $\{h,\ev{console}\}$.
Compute $\roots{\{k,\ev{unit}\}}$. Can evidence
$\{k\}\le_c\{\ev{log}\}$ be sound over this typed store when
$\kappa_{\ell}$ and $\kappa_c$ are distinct? Explain using root inclusion.}

\lesson{Immediate uses and result captures}\label{sec:latent}
\focus{A call's immediate uses can differ from the captures of its returned closure.}

To distinguish immediate uses from result captures, consider a factory
that returns a closure associated with another capability.
Figure~\ref{fig:factory} presents the factory and two clients in
simplified target notation, including their closing and avoidance
certificates. Calling $\ev{make}$ invokes $\ev{log}$ and constructs the
inner closure; its call to $\ev{console}$ remains latent until the
returned closure is invoked.


\begin{figure}[!hbp]
\begin{examplepanel}{tgtcol}{Factory and certificates in simplified target notation}
\small
\[
\tgt{\begin{aligned}
 v_{\mathrm{in}}
   &=\lambda^{\{\ev{console}\}}(v:\ty{Unit}).\
      \ev{console}\ v\ \langle g_{\mathrm{in}}\rangle,\\
 v_{\mathrm{make}}
   &=\lambda^{\{\ev{log}\}}(u:\ty{Unit}).\\[-2pt]
   &\quad
      \bigl(\ev{let}\ r=\ev{log}\ u\ \ev{in}\
         v_{\mathrm{in}}\ \langle\varnothing;q_0\rangle\bigr)
       \langle g_{\mathrm{make}}\rangle .
\end{aligned}}
\]
The symbol \(v_{\mathrm{in}}\) abbreviates the lambda literal itself,
with empty immediate uses. Its certificates and those of the factory are:
\[
\tgt{\begin{aligned}
 g_{\mathrm{in}}&=\ev{elem}(\{\ev{console},v\},\{\ev{console},v\}),\\
 q_0&=\ev{elem}(\varnothing,\varnothing),\\
 g_{\mathrm{make}}&=\ev{elem}(\{\ev{log},u\},\{\ev{log},u\}).
\end{aligned}}
\]
In a let annotation \(\langle U;q\rangle\), \(U\) bounds the continuation
body's uses and \(q\) proves the avoidance obligation.
The named certificates are checked under their displayed binders:
\(g_{\mathrm{in}}\) under \(v\), \(g_{\mathrm{make}}\) under \(u\),
and \(q_0\) under \(r\).
\end{examplepanel}
\begin{examplepanel}{tgtcol}{Clients of the stored factory \(\ev{make}\)}
\small
\[
\tgt{\begin{aligned}
 P_{\mathrm{call}}
   &=\ev{let}\ f=\ev{make}\ \ev{unit}\ \ev{in}\
      f\ \ev{unit}\ \langle U;q_{\mathrm{call}}\rangle,\\
 P_{\mathrm{drop}}
   &=\ev{let}\ f=\ev{make}\ \ev{unit}\ \ev{in}\
      \ev{unit}\ \langle\varnothing;q_{\mathrm{drop}}\rangle,\\
 U&=\{\ev{console},\ev{log}\},\\
 q_{\mathrm{call}}&:\{f,\ev{unit}\}\le_c U,\qquad
 q_{\mathrm{drop}}=\ev{capvar}(\ev{unit}):\{\ev{unit}\}\le_c\varnothing.
\end{aligned}}
\]
\end{examplepanel}
\begin{examplepanel}{runcol}{Initial roots of the whole state}
\[
\run{P_{\mathrm{call}}:\{\kappa_\ell,\kappa_c\},
      \qquad P_{\mathrm{drop}}:\{\kappa_\ell\}.}
\]
\end{examplepanel}
\caption{A factory with distinct immediate uses and result captures.
Both clients call \(\mathsf{make}\), which invokes \(\mathsf{log}\) and
returns a closure that invokes \(\mathsf{console}\);
\(P_{\mathrm{call}}\) then invokes that result.
The text constructs \(q_{\mathrm{call}}\) from two branches. This named
presentation follows the C1/C6 examples, with certificates supplied
directly as target evidence. Source translation may produce different
evidence syntax.}
\label{fig:factory}
\end{figure}



The outer closure can therefore have capture annotation $\{\ev{log}\}$,
while its result type has annotation $\{\ev{console}\}$:
\[
 \ev{make}:
 \bigl(\ty{Unit}\to
       (\ty{Unit}\to\ty{Unit})^{\{\ev{console}\}}\bigr)^{\{\ev{log}\}}.
\]
This typing relies on the empty immediate use set of a value expression,
even when its capture annotation is nonempty. The returned lambda's body
is checked with its own closing evidence and executes only when the
closure is called.

\remember{Capture annotations, immediate use sets, and syntactic free
variables are distinct. Effect safety bounds the capability roots
of variables inspected during execution; references reachable from a
returned closure or elsewhere in the heap require a separate property.}

\FloatBarrier
\subsection*{Continuation allowances}
Both clients in Figure~\ref{fig:factory} call $\ev{make}$, then differ in
how their let bodies use its result. The calling client invokes the
returned closure; the dropping client returns $\ev{unit}$. A bound based
only on $\ev{make}$'s annotation would omit the later console call in
$P_{\mathrm{call}}$, while charging every result capture to the initial
call would also charge $P_{\mathrm{drop}}$ for that unused capability.

The continuation supplies the additional information. While evaluating
$\ev{make}\ \ev{unit}$, the machine retains a frame for the pending
body. This frame declares an outer use set $U$ and carries evidence
that bounds the body's uses by $U$.

\needspace{12\baselineskip}
For $P_{\mathrm{call}}$, choose $U=\{\ev{console},\ev{log}\}$.
The continuation binds $f$ with annotation $\{\ev{console}\}$, allowing
the avoidance evidence to be constructed in two branches:
\[
 \begin{aligned}
 \{f\}&\xrightarrow{\ev{capvar}(f)}\{\ev{console}\}
          \xrightarrow{\ev{elem}}U,\\
 \{\ev{unit}\}&\xrightarrow{\ev{capvar}(\ev{unit})}\varnothing
          \xrightarrow{\ev{elem}}U.
 \end{aligned}
\]
Writing $q_f$ and $q_u$ for the two branch expressions, the certificate
$q_{\mathrm{call}}=\ev{union}(q_f,q_u)$ proves
$\{f,\ev{unit}\}\le_c U$. Its outer endpoint $U$ avoids the local name
$f$ by using names that account for its captures.

The chosen overapproximation $U$ retains $\ev{log}$, making the decrease
visible when the calling frame is discharged. The body of
$P_{\mathrm{drop}}$ uses only $\ev{unit}$, so
$U=\varnothing$ suffices, with evidence
$q_{\mathrm{drop}}=\ev{capvar}(\ev{unit}):
\{\ev{unit}\}\le_c\varnothing$.

Inside the factory, the let body is the inner lambda literal. Its empty
use set requires only
$q_0=\ev{elem}(\varnothing,\varnothing)$.
The factory body's use set is consequently
$\{\ev{log},u\}\cup\varnothing=\{\ev{log},u\}$.
This equals the allowance justified by the closing certificate
$g_{\mathrm{make}}$.

\subsection*{Whole-state roots along execution}
Before the factory call, the current expression contributes
$\{\ev{make},\ev{unit}\}$. The pending body contributes its declared $U$.
The initial roots are therefore
\[
 \begin{aligned}
 P_{\mathrm{call}}&:\ \{\kappa_{\ell},\kappa_c\},\\
 P_{\mathrm{drop}}&:\ \{\kappa_{\ell}\}.
 \end{aligned}
\]
Figure~\ref{fig:factory-run} follows the calling client through execution.


\begin{figure}[!htbp]
\begin{examplepanel}{runcol}{Whole-state roots for the calling client}
\small
\renewcommand{\arraystretch}{1.3}
\begin{tabularx}{\linewidth}{@{}>{\raggedright\arraybackslash}Xll@{}}
\toprule
Stage & Whole-state roots & Inspection\\
\midrule
Call \(\mathsf{make}\); keep the calling frame
 & $\{\kappa_\ell,\kappa_c\}$ & $\mathsf{make}$\\
Run its call to \(\mathsf{log}\); keep that frame
 & $\{\kappa_\ell,\kappa_c\}$ & $\mathsf{log}$\\
Inner closure ready for allocation
 & $\{\kappa_\ell,\kappa_c\}$ & none\\
Call the newly allocated \(f\); the frame is discharged
 & $\{\kappa_c\}$ & $f$\\
Run \(\mathsf{console}\ \mathsf{unit}\)
 & $\{\kappa_c\}$ & $\mathsf{console}$\\
Return \(\mathsf{unit}\)
 & $\varnothing$ & none\\
\bottomrule
\end{tabularx}
\end{examplepanel}
\caption{Execution of \(P_{\mathrm{call}}\) from
Figure~\ref{fig:factory}, with administrative transitions grouped into
rows. Before allocation, the pending frame contributes
\(U=\{\mathsf{console},\mathsf{log}\}\).
After that frame is discharged, the call's use set resolves to the console
root alone. The dropping client \(P_{\mathrm{drop}}\) allocates the result
and returns \(\mathsf{unit}\) without invoking it.}
\label{fig:factory-run}
\end{figure}



The trace groups administrative transitions. Discharging the calling
frame replaces its overapproximation, which retains $\ev{log}$, with
the use set of the actual call.

The dropping client also produces and allocates the inner closure, then
returns $\ev{unit}$ without invoking it. The console root is absent from
the initial root set and never occurs among the roots of an inspected
variable.

\exercise{latent}{Suppose the calling client instead tried to declare
$U=\{\ev{log}\}$. Which branch of its avoidance proof would fail?
Why does the same objection not apply to the dropping client?}

\lesson{Capture prediction and its scope}\label{sec:capture-scope}
\focus{Effect safety concerns runtime variable inspection; scope escape and classification require further hypotheses.}

Capture prediction states that the roots of a typed state's use set do
not grow along execution, with renaming accounting for store extension
by allocation.
A second lemma establishes that the roots of each inspected variable
are included in the roots of the current use set. Together these imply
effect safety: a capability root absent from the roots of the initial
use set cannot be a root of any later inspected variable.

The observable events are variable inspections, such as reading the
function receiver of an application or the object receiver of a
projection. Relating these observations to operating-system effects
would require semantics for those effects. Unboxing, which opens a
package that hides captures, has a separate corresponding safety result.

\subsection*{Capture members and inconsistent bounds}
Capture members expose lower and upper bounds, paralleling type members.
The name $x.C$ selects a capture member, as $x.A$ selects a type member;
$\{x.C\}$ is a capture set containing that selection.
Under a hypothetical interface
\[
 x:\{C:\{\kappa_{\ell}\}\,..\,\varnothing\},
\]
reading and composing its bounds gives
\[
 \{\kappa_{\ell}\}\le_c\{x.C\}\le_c\varnothing.
\]
This evidence is derivable under the hypothetical assumption. If a
typed store realized the interface with rigid
$\kappa_{\ell}$, its semantic interpretation would imply
$\{\kappa_{\ell}\}\subseteq\varnothing$, a contradiction.
The argument follows the same pattern as inconsistent type bounds:
the derivation is conditional on an interface that no typed store can
realize.

\needspace{8\baselineskip}
\subsection*{Scope escape}
An effect bound constrains what a computation may inspect. Scope rules
add conditions on whether a locally introduced capability can be returned
or subsumed into an outer capability set. Establishing such conditions
requires an additional argument beyond effect safety.

The \texttt{CapturesCC} extension introduces scope roots, levels, and
fresh capabilities. Its no-escape results depend on explicit restrictions;
in particular, relevant source results restrict derivations that read
capture bounds. These results exclude particular callback subcapturing
judgments under stated assumptions about scopes and bound reads.
Extending them to general \texttt{withFile} programs would require a
corresponding argument for those programs.

\subsection*{Classified capture bounds}
The \texttt{Classifiers} extension labels capabilities with classifiers,
such as a class of control capabilities. A kind specifies the permitted
classifiers, and kinding evidence establishes that a capture set belongs
to that kind. Filtered capture sets describe the corresponding roots;
classified effect safety ensures that inspected roots satisfy the
permitted classification.

This extends the root-set argument with kinding evidence and its
interpretation while retaining the store, substitution, and prediction
reasoning. Operational results for exception handlers or thread
scheduling would additionally require models of those mechanisms.
